\section*{期中复习卷}
\subsection*{导数}
\begin{enumerate}
    \item 设函数 $y = f\left( x\right)$ 在 $x = {x}_{0}$ 处可导,且 ${f}^{\prime }\left( {x}_{0}\right)  = 1$ ,则 $\mathop{\lim }\limits_{{h \rightarrow  0}}\dfrac{f\left( {{x}_{0} + {3h}}\right)  - f\left( {x}_{0}\right) }{h} =$ \blkda{3}.
    \begin{jieda}
        根据导数的定义可知: $\mathop{\lim }\limits_{{h \rightarrow  0}}\dfrac{f\left( {{x}_{0} + {3h}}\right)  - f\left( {x}_{0}\right) }{h} = 3\mathop{\lim }\limits_{{h \rightarrow  0}}\dfrac{f\left( {{x}_{0} + {3h}}\right)  - f\left( {x}_{0}\right) }{3h} = 3{f}^{\prime }\left( {x}_{0}\right)  = 3$.
    \end{jieda}
    
    \item 给出定义: 设 ${f}^{\prime }\left( x\right)$ 是函数 $y = f\left( x\right)$ 的导函数, ${f}^{\prime \prime }\left( x\right)$ 是函数 $y = {f}^{\prime }\left( x\right)$ 的导函数,若方程 ${f}^{\prime \prime }\left( x\right)  = 0$ 有实数解 $x = {x}_{0}$ ,则称 $\left( {{x}_{0},f\left( {x}_{0}\right) }\right)$ 为函数 $y = f\left( x\right)$ 的 “拐点”. 经研究发现所有的三次函数 $f\left( x\right)  = a{x}^{3} + b{x}^{2} + {cx} + d\left( {a \neq  0}\right)$ 都有“拐点”,且该“拐点”也是函数 $y = f\left( x\right)$ 的图像的对称中心. 若函数 $f\left( x\right)  =  - {x}^{3} + 3{x}^{2}$ , 则 $f\left( \dfrac{1}{2021}\right)  + f\left( \dfrac{2}{2021}\right)  + f\left( \dfrac{3}{2021}\right)  + \cdots  + f\left( \dfrac{4040}{2021}\right)  + f\left( \dfrac{4041}{2021}\right)  =$\dotfill{(\kh{A})}
    \xx{8082}{-8082}{8084}{-8084}
    
    \item 一个球的半径以 2 厘米/秒的速度在增长，当此球的半径为 6 厘米时，其体积的变化率为\dotfill{(\kh{B})}
    \xx{$576\pi {cm}^{3}/$ 秒}{$288\pi {cm}^{3}/$ 秒}{$144\pi {cm}^{3}/$ 秒}{${96\pi } {cm}^{3}/$ 秒}
    
    \item 已知函数 $y = f\left( x\right)$ 及其导函数 $y = {f}^{\prime }\left( x\right)$ 的定义域均为$\mathbf{R}$,记 $g\left( x\right)  = {f}^{\prime }\left( x\right)$ . 若 $y = f\left( {\dfrac{3}{2} - {2x}}\right) \text{ 、 }y = g\left( {2 + x}\right)$ 均为偶函数,有下列两个结论:  \ding{172}  $f\left( {-1}\right)  = f\left( 4\right)$ ;  \ding{173}  $g\left( {-\dfrac{1}{2}}\right)  = 0$ ，则\dotfill{(\kh{C})}
    \xx{\ding{172} 对 \ding{173} 错}{\ding{172} 错 \ding{173} 对}{\ding{172} 对 \ding{173} 对}{\ding{172} 错 \ding{173} 错}
    \begin{jieda}
        $f(x)$关于$x = \dfrac{3}{2}$对称，故\ding{172}正确，$g(x)$关于$x = 2$对称

        因为两函数定义域均为$\mathbf{R}$，故$g(x)$有对称中心$\left(\dfrac{3}{2}, 0\right)$，故$g(x)$周期为$2$，故\ding{173}正确.
    \end{jieda}
    
    \item 求曲线 $y = \cos x$ 在 $x = \dfrac{\pi }{2}$ 处的切线方程是\blkda[3]{$y = -\left(x-\dfrac{\pi}{2}\right)$}.
    \begin{jieda}
        设$f(x) = cosx$，$f'(x) = -sinx$，$f'\left(\dfrac{\pi}{2}\right) = -1$，故切线为$y = -\left(x-\dfrac{\pi}{2}\right)$
    \end{jieda}

    \item 求过点 $\left( {0, - 1}\right)$ 且与曲线 $y = 2{x}^{2}$ 相切的直线的方程是\blkda[3]{$y = \pm 2\sqrt{2}x-1$}.
    \begin{jieda}
        设切点为$(m, 2m^2)$，该处切线为$y = 4m(x-m)+2m^2$，令其过点 $\left( {0, - 1}\right)$，得到$m = \pm \dfrac{\sqrt{2}}{2}$，故切线为$y = \pm 2\sqrt{2}x-1$
    \end{jieda}
    
    \item 直线 $l : y = {kx} + b$ 是曲线 $y = \ln x$ 和 $y = {{e}}^{x - 2}$ 的公切线，则 $k + b =$ \blkda{$\dfrac{1}{e}$或0}.
    \begin{jieda}
        对于 $y = \ln x$ ,设切点为 $\left( {{x}_{1},\ln {x}_{1}}\right)$ ,求导得 ${y}^{\prime } = \dfrac{1}{x}$，则在该点处的斜率为 $k = \dfrac{1}{{x}_{1}}$，则切线方程为: $y - \ln {x}_{1} = \dfrac{1}{{x}_{1}}\left( {x - {x}_{1}}\right)$ ,即 $y = \dfrac{1}{{x}_{1}}x + \ln {x}_{1} - 1$
        
        对于 $y = {{e}}^{x - 2}$ ,设切点为 $\left( {{x}_{2},{{e}}^{{x}_{2} - 2}}\right)$ ,求导得 ${y}^{\prime } = {{e}}^{x - 2}$，则在该点处的斜率为 $k = {{e}}^{{x}_{2} - 2}$，则切线方程为: $y - {{e}}^{{x}_{2} - 2} = {{e}}^{{x}_{2} - 2}\left( {x - {x}_{2}}\right)$ ,即 $y = {{e}}^{{x}_{2} - 2}x + \left( {1 - {x}_{2}}\right) {{e}}^{{x}_{2} - 2}$
        
        因为 $l : y = {kx} + b$ 是公切线，所以 $\left\{  \begin{array}{l} \dfrac{1}{{x}_{1}} = {{e}}^{{x}_{2} - 2} \\  \ln {x}_{1} - 1 = \left( {1 - {x}_{2}}\right) {{e}}^{{x}_{2} - 2} \end{array}\right.$ ,即 ${x}_{1} = {{e}}^{2 - {x}_{2}}$，所以 $\ln {{e}}^{2 - {x}_{2}} - 1 = \left( {1 - {x}_{2}}\right) {{e}}^{{x}_{2} - 2}$ ,即 $1 - {x}_{2} = \left( {1 - {x}_{2}}\right) {{e}}^{{x}_{2} - 2}$，所以 $\left( {1 - {x}_{2}}\right) \left( {{{e}}^{{x}_{2} - 2} - 1}\right)  = 0$，即 $1 - {x}_{2} = 0$ 或 ${x}_{2} - 2 = 0$ ,解得 ${x}_{2} = 1$ 或 ${x}_{2} = 2$
        
        当 ${x}_{2} = 1$ 时,此时 $k = {{e}}^{{x}_{2} - 2} = \dfrac{1}{{e}},b = \left( {1 - {x}_{2}}\right) {{e}}^{{x}_{2} - 2} = 0$ ,所以 $k + b = \dfrac{1}{{e}}$
        
        当 ${x}_{2} = 2$ 时,此时 $k = {{e}}^{{x}_{2} - 2} = 1,b = \left( {1 - {x}_{2}}\right) {{e}}^{{x}_{2} - 2} =  - 1$ ,所以 $k + b = 0$
        
        所以 $k + b = 0$ 或 $\dfrac{1}{{e}}$
    \end{jieda}

    \item 函数 $f\left( x\right)  = {x}^{2} - {x}^{3}$ 的严格递增区间为\blkda{$\left(0, \dfrac{2}{3}\right)$}，严格递减区间为\blkda[3]{$(-\infty, 0), \left(\dfrac{2}{3}, +\infty \right)$}.
    
    \item 已知函数 $f\left( x\right)  = \dfrac{1}{3}{a}^{2}{x}^{3} - \dfrac{5}{2}a{x}^{2} + {6x}$ 的极小值点为 1，则 $a =$ \blkda{3}.
    \begin{jieda}
        易知 ${f}^{\prime }\left( x\right)  = {a}^{2}{x}^{2} - {5ax} + 6 = \left( {{ax} - 2}\right) \left( {{ax} - 3}\right)$，由题意可知 $x = 1$ 是其一个变号零点，即 ${f}^{\prime }\left( 1\right)  = \left( {a - 2}\right) \left( {a - 3}\right)  = 0 \Longrightarrow  a = 2$ 或 $a = 3$
        \begin{dingautolist}{172}
            \item 当 $a = 2$ 时, ${f}^{\prime }\left( x\right)  = 2\left( {x - 1}\right) \left( {{2x} - 3}\right)$，此时 $x \in  \left( {1,\dfrac{3}{2}}\right)$ 时, ${f}^{\prime }\left( x\right)  < 0$ ,即 $f\left( x\right)$ 单调递减，$x \in  \left( {-\infty ,1}\right)$ , ${f}^{\prime }\left( x\right)  > 0$ ,即 $f\left( x\right)$ 单调递增,故在 $x = 1$ 处取得极大值,不符合题意
            \item 当 $a = 3$ 时, ${f}^{\prime }\left( x\right)  = 3\left( {{3x} - 2}\right) \left( {x - 1}\right)$，此时 $x \in  \left( {\dfrac{2}{3},1}\right)$ 时, ${f}^{\prime }\left( x\right)  < 0$ ,即 $f\left( x\right)$ 单调递减，$x \in  \left( {1, + \infty }\right)$ , ${f}^{\prime }\left( x\right)  > 0$ ,即 $f\left( x\right)$ 单调递增,故在 $x = 1$ 处取得极小值,符合题意
        \end{dingautolist}
        综上 $a = 3$ .
    \end{jieda}
    
    \item 函数 $f\left( x\right)  = {x}^{3} + a{x}^{2} + {bx} + {a}^{2}$ 在 $x = 1$ 时有极值 10, 那么 $a$ 、 $b$ 的值为\blkda[4]{$a = -4, b = 11$}.
    \begin{jieda}
        因为函数 $f\left( x\right)  = {x}^{3} - a{x}^{2} - {bx} + {a}^{2}$ ,
        所以 ${f}^{\prime }\left( x\right)  = 3{x}^{2} - {2ax} - b$ ,
        因为函数 $f\left( x\right)  = {x}^{3} - a{x}^{2} - {bx} + {a}^{2}$ 在 $x = 1$ 时有极值 10,
        所以 $\left\{  \begin{array}{l} 1 - a - b + {a}^{2} = {10} \\  3 - {2a} - b = 0 \end{array}\right.$ ,
        即 $\left\{  \begin{array}{l}  - a - b + a^2 = 9 \\  {2a} + b = 3 \end{array}\right.$ ,
        解得 $\left\{  {\begin{array}{l} a =  - 4 \\  b = {11} \end{array}\text{ 或 }\left\{  {\begin{array}{l} a = 3 \\  b =  - 3 \end{array},}\right. }\right.$
        当 $\left\{  \begin{array}{l} a =  - 4 \\  b = {11} \end{array}\right.$ 时, ${f}^{\prime }\left( x\right)  = 3{x}^{2} + {8x} - {11} = \left( {{3x} + {11}}\right) \left( {x - 1}\right)$
        当 $x <  - \dfrac{11}{3}$ 或 $x > 1$ 时, ${f}^{\prime }\left( x\right)  > 0$ ,当 $- \dfrac{11}{3} < x < 1$ 时, ${f}^{\prime }\left( x\right)  < 0$ .
        所以 $f\left( x\right)$ 在 $x = 1$ 处取的极小值,符合题意.
        当 $\left\{  \begin{array}{l} a = 3 \\  b =  - 3 \end{array}\right.$ 时, ${f}^{\prime }\left( x\right)  = 3{x}^{2} - {6x} + 3 = 3{\left( x - 1\right) }^{2} \geq  0$
        所以 $f\left( x\right)$ 在 $x = 1$ 处无极值.
    \end{jieda}


    \item 已知 $f\left( x\right)  = {{e}}^{x} - {ax} - 1$ .
    \begin{enumerate}
        \item 当 $a = 2$ 时，讨论 $f\left( x\right)$ 的单调区间
        \item 若 $f\left( x\right)$ 在定义域 ${R}$ 内严格单调递增，求 $a$ 的取值范围
    \end{enumerate}
    \begin{jieda}
        \begin{enumerate}
            \item 当 $a = 2$ 时, ${f}^{\prime }\left( x\right)  = {e}^{x} - 2$ .
            令 ${f}^{\prime }\left( x\right)  > 0$ ,即 ${e}^{x} - 2 > 0$ 解得： $x > \ln 2$ ；
            令 ${f}^{\prime }\left( x\right)  < 0$ ,即 ${e}^{x} - 2 < 0$ 解得： $x < \ln 2$ ；
            $\therefore$ 当 $a = 2$ 时,函数 $f\left( x\right)$ 的单调增区间是 $\left( {\ln 2, + \infty }\right)$ ,递减区间为 $\left( {-\infty ,\ln 2}\right)$ .
            \item $\because f\left( x\right)  = {e}^{x} - {ax} - 1,x \in  \mathbf{R}\therefore {f}^{\prime }\left( x\right)  = {e}^{x} - a$
            $\because f\left( x\right)$ 在 ${R}$ 上单调递增,即 ${f}^{\prime }\left( x\right)  = {e}^{x} - a \geq  0$ 恒成立,
            $\because x \in  {R}$ 时 ${e}^{x} \in  \left( {0, + \infty }\right)$
            $\therefore a \leq  0$ ,即 $a$ 的取值范围为 $( - \infty ,0\rbrack$ .
        \end{enumerate}
    \end{jieda}
    
    \item 已知 $f\left( x\right)  = a{x}^{2} + x - \ln x$ .
    \begin{enumerate}
        \item 若 $a = 1$ ，求函数 $y = f\left( x\right)$ 的极值
        \item 若 $y = f\left( x\right)$ 存在单调递增区间，求实数 $a$ 的取值范围
    \end{enumerate}
    \begin{jieda}
        \begin{enumerate}
            \item $a = 1$ 时, $f\left( x\right)  = {x}^{2} + x - \ln x$ 定义域为 $\left( {0, + \infty }\right)$ ,
            ${f}^{\prime }\left( x\right)  = {2x} + 1 - \dfrac{1}{x} = \dfrac{2{x}^{2} + x - 1}{x} = \dfrac{\left( {{2x} - 1}\right) \left( {x + 1}\right) }{x},x > 0$
            令 ${f}^{\prime }\left( x\right)  = 0$ ,解得 $x = \dfrac{1}{2}$ .
            当 $x \in  \left( {0,\dfrac{1}{2}}\right)$ 时, ${f}^{\prime }\left( x\right)  < 0,f\left( x\right)$ 单调递减;
            当 $x \in  \left( {\dfrac{1}{2}, + \infty }\right)$ 时, ${f}^{\prime }\left( x\right)  > 0,f\left( x\right)$ 单调递增.
            所以在 $x = \dfrac{1}{2}$ 时,函数 $f\left( x\right)$ 取得极小值 $f\left( \dfrac{1}{2}\right)  = \dfrac{3}{4} + \ln 2;f\left( x\right)$ 无极大值.
            \item ${f}^{\prime }\left( x\right)  = {2ax} + 1 - \dfrac{1}{x} = \dfrac{{2a}{x}^{2} + x - 1}{x},\left( {x > 0}\right)$
            因为 $y = f\left( x\right)$ 存在单调递增区间,所以 ${2a}{x}^{2} + x - 1 > 0$ 在 $\left( {0, + \infty }\right)$ 上有解.
            由 ${2a}{x}^{2} + x - 1 > 0,x \in  \left( {0, + \infty }\right)$ 得 ${2a} > \dfrac{1}{{x}^{2}} - \dfrac{1}{x} = {\left( \dfrac{1}{x} - \dfrac{1}{2}\right) }^{2} - \dfrac{1}{4}$ ,
            $\because {\left( \dfrac{1}{x} - \dfrac{1}{2}\right) }^{2} - \dfrac{1}{4} \geq   - \dfrac{1}{4}$ ,当且仅当 $x = 2$ 时取等号.
            $\therefore {2a} >  - \dfrac{1}{4}$ ,即 $a >  - \dfrac{1}{8}$ ,
            所以实数 $a$ 的取值范围是 $\left( {-\dfrac{1}{8}, + \infty }\right)$ .
        \end{enumerate}
    \end{jieda}
    \item 已知函数 $f\left( x\right)  = 3\ln x - x + \dfrac{a}{x}$ .
    \begin{enumerate}
        \item 若 $a = 2$ ，求 $f\left( x\right)$ 的单调区间
        \item 若 $f\left( x\right)$ 没有极值点，求 $a$ 的取值范围
    \end{enumerate}
    \begin{jieda}
        \begin{enumerate}
            \item 若 $a = 2$ ,则 $f\left( x\right)  = 3\ln x - x + \dfrac{2}{x}$ ,函数定义域为 $\left( {0, + \infty }\right)$ ,
            ${f}^{\prime }\left( x\right)  = \dfrac{3}{x} - 1 - \dfrac{2}{{x}^{2}} =  - \dfrac{\left( {x - 1}\right) \left( {x - 2}\right) }{{x}^{2}}.$
            当 $0 < x < 1$ 时, ${f}^{\prime }\left( x\right)  < 0$ ；
            当 $1 < x < 2$ 时, ${f}^{\prime }\left( x\right)  > 0$ ；
            当 $x > 2$ 时, ${f}^{\prime }\left( x\right)  < 0$ ,
            故 $f\left( x\right)$ 的单调递减区间是$(0,1)$和 $\left( {2, + \infty }\right)$ ,单调递增区间是$(1,2)$.
            \item ${f}^{\prime }\left( x\right)  = \dfrac{3}{x} - 1 - \dfrac{a}{{x}^{2}} =  - \dfrac{{x}^{2} - {3x} + a}{{x}^{2}}$ ,
            函数 $h\left( x\right)  = {x}^{2} - {3x} + a$ ,当 $\Delta  = 9 - {4a} \leq  0$ ,即 $a \geq  \dfrac{9}{4}$ 时, $h\left( x\right)  \geq  0$ 恒成立，则有 ${f}^{\prime }\left( x\right)  \leq  0,f\left( x\right)$ 单调递减，此时 $f\left( x\right)$ 没有极值点,符合题意.
            当 $a < \dfrac{9}{4}$ 时,方程 ${x}^{2} - {3x} + a = 0$ 有两个实数根 ${x}_{1},{x}_{2}$ ,不妨设 ${x}_{1} < {x}_{2}$ ,
            则 ${x}_{1} + {x}_{2} = 3 > 0,{x}_{1}{x}_{2} = a$ .
            当 $0 < a < \dfrac{9}{4}$ 时, $0 < {x}_{1} < {x}_{2}$ ,此时 $f\left( x\right)$ 在区间 $\left( {0,{x}_{1}}\right) ,\left( {{x}_{2}, + \infty }\right)$ 上单调递减,
            在区间 $\left( {{x}_{1},{x}_{2}}\right)$ 上单调递增,所以 ${x}_{1}$ 是 $f\left( x\right)$ 的极小值点, ${x}_{2}$ 是 $f\left( x\right)$ 的极大值点,不符合题意;
            同理可知,当 $a \leq  0$ 时, $f\left( x\right)$ 在区间 $\left( {0,{x}_{2}}\right)$ 上单调递增, $\left( {{x}_{2}, + \infty }\right)$ 上单调递减, ${x}_{2}$ 是 $f\left( x\right)$ 的
            极大值点, 不符合题意.
            综上, $a$ 的取值范围是 $\left\lbrack  {\dfrac{9}{4}, + \infty }\right)$
        \end{enumerate}
    \end{jieda}

    
    \item 已知函数 $f\left( x\right)  = \left\{  \begin{array}{l} {{e}}^{x},x > 0 \\  {x}^{3} - {3x} + a,x \leq  0 \end{array}\right.$ 的值域为 $\mathbf{R}$ ，则实数 $a$ 的取值范围为\blkda{$[-1, +\infty)$}.
    \begin{jieda}
        当 $x > 0$ 时, $f\left( x\right)  = {e}^{x}$ 的值域为 $\left( {1, + \infty }\right)$

        故应有当 $x \leq  0$ 时, $( - \infty ,1\rbrack$ 是 $f\left( x\right)  = {x}^{3} - {3x} + a$ 的值域的子集，又 ${f}^{\prime }\left( x\right)  = 3{x}^{2} - 3 = 3\left( {x + 1}\right) \left( {x - 1}\right)$ ,令 ${f}^{\prime }\left( x\right)  = 0,\therefore x =  - 1$ 或 $x = 1$ (舍去)
        
        当 $x <  - 1$ 时, ${f}^{\prime }\left( x\right)  > 0$ , $f\left( x\right)$ 单调递增；当 $- 1 < x \leq  0$ 时, ${f}^{\prime }\left( x\right)  < 0,f\left( x\right)$ 单调递减
        
        当 $x =  - 1$ 时, $f\left( x\right)$ 取得极大值 $f\left( {-1}\right)  = 2 + a$，故 $f\left( x\right)  = {x}^{3} - {3x} + a\left( {x \leq  0}\right)$ 的值域为 $( - \infty ,2 + a\rbrack$
        
        故$2 + a \geq  1$，故$a \geq   - 1$ .
    \end{jieda}

    
    \item 若函数 $f\left( x\right)  =  -{x}^{3} + \dfrac{3}{2}{x}^{2}$ 在区间 $\left( {{2a} - 5,a + 1}\right)$ 内有最小值，则实数 $a$ 的取值范围是\blkda{$\left(-1, \dfrac{1}{2}\right]$}.
    \begin{jieda}
        由题可知: ${f}^{\prime }\left( x\right)  =  - 3{x}^{2} + {3x} =  - {3x}\left( {x - 1}\right)$，令 ${f}^{\prime }\left( x\right)  > 0$，则 $0 < x < 1$ ; 令 ${f}^{\prime }\left( x\right)  < 0$ ,则 $x < 0$ 或 $x > 1$，所以函数 $f\left( x\right)$ 在$(0,1)$单调递增,在 $\left( {-\infty ,0}\right) ,\left( {1, + \infty }\right)$ 单调递减.
        
        极小值为 $f\left( 0\right)  = 0$ ,令 $f\left( x\right)  = 0$ ,所以 $x = 0$ 或 $\dfrac{3}{2}$，又函数 $f\left( x\right)  =  - {x}^{3} + \dfrac{3}{2}{x}^{2}$ 在区间 $\left( {{2a} - 5,a + 1}\right)$ 内有最小值，所以 $\left\{  \begin{array}{l} {2a} - 5 < 0 \\  0 < a + 1 \leq  \dfrac{3}{2} \end{array}\right.  \Longrightarrow   - 1 < a \leq  \dfrac{1}{2}$.
    \end{jieda}
    
    \item 已知函数 $f\left( x\right)  = \ln x - \dfrac{1}{2}{x}^{2}$ . 
    \begin{enumerate}
        \item 求函数 $f\left( x\right)$ 在 $\left\lbrack  {\dfrac{1}{2},4}\right\rbrack$ 上的最大值和最小值
        \item 若不等式 $f\left( x\right)  > \left( {2 - a}\right) {x}^{2}$ 有解,求实数 $a$ 的取值范围
    \end{enumerate}
    \begin{jieda}
        \begin{enumerate}
            \item 因为函数 $f\left( x\right)  = \ln x - \dfrac{1}{2}{x}^{2}$，所以 ${f}^{\prime }\left( x\right)  = \dfrac{1}{x} - x = \dfrac{1 - {x}^{2}}{x}\left( {x > 0}\right)$ .令 ${f}^{\prime }\left( x\right)  = 0$ ,则 $x = 1$ 或 $x =  - 1$ (舍去).
            
            当 $\dfrac{1}{2} < x < 1$ 时, ${f}^{\prime }\left( x\right)  > 0$ ,当 $1 < x < 4$ 时, ${f}^{\prime }\left( x\right)  < 0$
            
            所以函数 $f\left( x\right)$ 在 $\left( {\dfrac{1}{2},1}\right)$ 上单调递增,在$(1,4)$上单调递减，所以当 $x = 1$ 时, $f\left( x\right)$ 取得最大值. 最大值为 $f\left( 1\right)  =  - \dfrac{1}{2}$，又 $f\left( \dfrac{1}{2}\right)  = \ln \dfrac{1}{2} - \dfrac{1}{8} =  - \ln 2 - \dfrac{1}{8},f\left( 4\right)  = 2\ln 2 - 8$
            
            $f\left( \dfrac{1}{2}\right)  - f\left( 4\right)  =  - 3\ln 2 - \dfrac{1}{8} + 8 >  - 3 - \dfrac{1}{8} + 8 > 0.$
            
            所以 $f\left( \dfrac{1}{2}\right)  > f\left( 4\right)$ . 所以当 $x = 4$ 时, $f\left( x\right)$ 取得最小值. 最小值为 $f\left( 4\right)  = 2\ln 2 - 8$
            
            故 $f\left( x\right)$ 在 $\left\lbrack  {\dfrac{1}{2},4}\right\rbrack$ 上的最大值为 $- \dfrac{1}{2}$ . 最小值为 $2\ln 2 - 8$.
            
            \item 易知 $f\left( x\right)$ 的定义域为 $\left( {0, + \infty }\right)$ .故不等式 $f\left( x\right)  > \left( {2 - a}\right) {x}^{2}$ 可化为 $2 - a < \dfrac{\ln x}{{x}^{2}} - \dfrac{1}{2}$.
            
            记 $g\left( x\right)  = \dfrac{\ln x}{{x}^{2}} - \dfrac{1}{2}$ . 则原不等式有解可转化为 $2 - a < g{\left( x\right) }_{\max }$，易得 ${g}^{\prime }\left( x\right)  = \dfrac{1 - 2\ln x}{{x}^{3}},x \in  \left( {0,\sqrt{{e}}}\right)$ 时, ${g}^{\prime }\left( x\right)  > 0,x \in  \left( {\sqrt{{e}}, + \infty }\right)$ 时, ${g}^{\prime }\left( x\right)  < 0$，所以函数 $g\left( x\right)$ 在 $\left( {0,\sqrt{{e}}}\right)$ 上单调递增,在 $\left( {\sqrt{{e}}, + \infty }\right)$ 上单调递减，故 $g{\left( x\right) }_{\max } = g\left( \sqrt{{e}}\right)  = \dfrac{1}{2{e}} - \dfrac{1}{2}$ ,所以 $2 - a < \dfrac{1}{2{e}} - \dfrac{1}{2}$
            
            解得 $a > \dfrac{5}{2} - \dfrac{1}{2{e}}$.
        \end{enumerate}
    \end{jieda}

    \item 已知函数 $f\left( x\right)  = \dfrac{a\ln x}{x} + x - 1$ .
    \begin{enumerate}
        \item 当 $a = 1$ 时，求函数 $g\left( x\right)  = f\left( x\right)  - x$ 极值
        \item 若对任意 $x \in  \lbrack 1, + \infty ),f\left( x\right)  \geq  a + 1$ 恒成立,求实数 $a$ 的取值范围
    \end{enumerate}
    \begin{jieda}
        \begin{enumerate}
            \item 函数 $f\left( x\right)  = \dfrac{a\ln x}{x} + x - 1$ 的定义域为 $\left( {0, + \infty }\right)$ . 当 $a = 1$ 时, $g\left( x\right)  = f\left( x\right)  - x = \frac{\ln x}{x} - 1$.
            求导得 ${g}^{\prime }\left( x\right)  = \dfrac{1 - \ln x}{{x}^{2}}$ . 由 ${g}^{\prime }\left( x\right)  = 0$ . 得 $x = {e}$ . 由 ${g}^{\prime }\left( x\right)  > 0$ . 得 $0 < x < {e}$ . 由 ${g}^{\prime }\left( x\right)  < 0$ ,得 $x > {e}$ .
            因此 $g\left( x\right)$ 在$(0,e)$上单调递增. 在 $\left( {{e}, + \infty }\right)$ 上单调递减.
            所以 $g\left( x\right)$ 在 $x = {e}$ 处取得极大值 $g\left( {e}\right)  = \dfrac{1}{{e}} - 1$ ,无极小值.
            \item 函数 $f\left( x\right)  = \dfrac{a\ln x}{x} + x - 1$ , $f\left( x\right)  \geq  a + 1 \Longleftrightarrow  a\left( {\ln x - x}\right)  \geq  {2x} - {x}^{2}$ , $x \in  \lbrack 1, + \infty )$ .
            设 $m\left( x\right)  = \ln x - x.x \in  \lbrack 1, + \infty )$ . 求导得 ${m}^{\prime }\left( x\right)  = \dfrac{1 - x}{x} \leq  0$ . 函数 $m\left( x\right)$ 在 $\lbrack 1, + \infty )$ 上单调递减.
            则 $m\left( x\right)  \leq  m\left( 1\right)  =  - 1 < 0$ . 即 $\ln x < x$ . 因此 $a \leq  \dfrac{{x}^{2} - {2x}}{x - \ln x}$ .
            令 $\varphi \left( x\right)  = \dfrac{{x}^{2} - {2x}}{x - \ln x},x \in  \lbrack 1, + \infty )$ . 求导得 ${\varphi }^{\prime }\left( x\right)  = \dfrac{\left( {x - 1}\right) \left( {x + 2 - 2\ln x}\right) }{{\left( x - \ln x\right) }^{2}}$ .
            令 $h\left( x\right)  = x + 2 - 2\ln x,x \in  \lbrack 1, + \infty )$ ,求导得 ${h}^{\prime }\left( x\right)  = 1 - \dfrac{2}{x}$ . 当 $1 \leq  x < 2$ 时, ${h}^{\prime }\left( x\right)  < 0$ ,
            当 $x > 2$ 时, ${h}^{\prime }\left( x\right)  > 0$ ,即 $h\left( x\right)$ 在(1,2)上单调递减,在 $\left( {2, + \infty }\right)$ 上单调递增.
            则 $h{\left( x\right) }_{\min } = h\left( 2\right)  = 4 - 2\ln 2 > 0$ ,即 ${\varphi }^{\prime }\left( x\right)  \geq  0$ ,因此函数 $\varphi \left( x\right)$ 在 $\lbrack 1, + \infty )$ 上是增函数.
            $\varphi {\left( x\right) }_{\min } = \varphi \left( 1\right)  =  - 1$
            所以 $a \leq   - 1$ ,即实数 $a$ 的取值范围为 $( - \infty , - 1\rbrack$ .
        \end{enumerate}
    \end{jieda}
    
        
        \item 已知函数$f(x) = \dfrac{1}{2}ax^2-(2a+1)x + 2lnx$（$a \in \mathbf{R}$）
        \begin{enumerate}
            \item 求$y = f(x)$单调区间
            \item 设$g(x) = x^2-2x$，若对任意的$x_1 \in (0, 2]$，均存在$x_2 \in (0, 2]$，使得$f(x_1) < g(x_2)$，求$a$的取值范围.
        \end{enumerate}
        \begin{jieda}
            \begin{enumerate}
                \item $f'(x) = ax-(2a+1) + \dfrac{2}{x} = \dfrac{(ax-1)(x-2)}{x} (x > 0)$
                \begin{dingautolist}{172}
                    \item 若$a < 0$，则$f(x)$在$(0, 2)$内严增，在$(2, +\infty)$严减
                    \item 若$a = 0$，则$f(x)$在$(0, 2)$内严增，在$(2, +\infty)$严减
                    \item 若$a \in \left(0, \dfrac{1}{2}\right)$，则$f(x)$在$(0, 2)$严增，在$\left(2, \dfrac{1}{a}\right)$严减，在$\left(\dfrac{1}{a}, +\infty\right)$严增.
                    \item 若$a = \dfrac{1}{2}$，则$f(x)$在$(0, +\infty)$严增
                    \item 若$a \in \left(\dfrac{1}{2}, +\infty\right)$，则$f(x)$在$\left(0, \dfrac{1}{a}\right)$严增，在$\left(\dfrac{1}{a}, 2\right)$严减，在$\left(2, +\infty\right)$严增.
                \end{dingautolist}
                \item 由题意可知$g(x)$在$(0, 2]$上的最大值大于$f(x)$在$(0, 2]$上的最大值. $g(x)$在$(0, 2]$上的最大值为$0$. 
                \begin{dingautolist}{172}
                    \item 若$a \leq \dfrac{1}{2}$，则$f(x) \leq f(2) = -2-2a+2ln2$，当$a \in \left(ln2-1, \dfrac{1}{2}\right]$时满足要求
                    \item 若$a \in \left(\dfrac{1}{2}, +\infty\right)$，则$f(x) \leq f\left(\dfrac{1}{a}\right) = -\dfrac{1}{2a} - 2-2lna$ \\
                    设$h(a) = -\dfrac{1}{2a} - 2-2lna$，$h'(a) = \dfrac{1}{2a^2} - \dfrac{2}{a}$，$h'(a) < 0$在$a \in \left(\dfrac{1}{2}, +\infty\right)$恒成立，故$h(a) < h\left(\dfrac{1}{2}\right) = -3+2ln2 < 0$，因此满足要求
                \end{dingautolist}
                综上，$a \in (ln2-1, +\infty)$
            \end{enumerate}
        \end{jieda}

        
        \item 已知定义在 $\left\lbrack  {-3,3}\right\rbrack$ 上的函数 $f\left( x\right)  = {{e}}^{x} - {{e}}^{-x} - {2x} + 1$ ，若 $f\left( {m}^{2}\right)  +  f\left( {m - 2}\right)  \leq  2$ ，则 $m$ 的取值范围是\blkda{$\left\lbrack  {-1,1}\right\rbrack$}

        \begin{jieda}
            记 $g\left( x\right)  = {{e}}^{x} - {{e}}^{-x} - {2x},x \in  \left\lbrack  {-3,3}\right\rbrack$ ,则 $g\left( {-x}\right)  = {{e}}^{-x} - {{e}}^{x} + {2x} =  - g\left( x\right)$，故 $g\left( x\right)$ 为 $\left\lbrack  {-3,3}\right\rbrack$ 的奇函数
            
            因为${g}^{\prime }\left( x\right)  = {{e}}^{x} + {{e}}^{-x} - 2 \geq  2\sqrt{{{e}}^{x}{{e}}^{-x}} - 2 = 0$，所以$g\left( x\right)$ 为 $\left\lbrack  {-3,3}\right\rbrack$ 上的单调递增函数
            
            因为 $f\left( x\right)  = g\left( x\right)  + 1$，由 $f\left( {m}^{2}\right)  + f\left( {m - 2}\right)  \leq  2$ 可得 $g\left( {m}^{2}\right)  + g\left( {m - 2}\right)  + 2 \leq  2$ ,进而 $g\left( {m}^{2}\right)  \leq  g\left( {2 - m}\right)$ ,
            
            故 $- 3 \leq  {m}^{2} \leq  2 - m \leq  3$ ，解得 $- 1 \leq  m \leq  1$
        \end{jieda}
        
        \item 定义在 $\left( {0, + \infty }\right)$ 上的函数 $f\left( x\right)$ 的导函数为 ${f}^{\prime }\left( x\right)$ ,若 $x{f}^{\prime }\left( x\right)  - f\left( x\right)  < 0$ ,且 $f\left( 3\right)  = 0$ ,则不等式 $\left( {x - 2}\right) f\left( x\right)  < 0$ 的解集为\blkda[3]{$\left( {0,2}\right)  \cup  \left( {3, + \infty }\right)$}
        \begin{jieda}
            依题意,令 $g\left( x\right)  = \dfrac{f\left( x\right) }{x}$ ,求导得 ${g}^{\prime }\left( x\right)  = \dfrac{x{f}^{\prime }\left( x\right)  - f\left( x\right) }{{x}^{2}} < 0$ ,则 $g\left( x\right)$ 在 $\left( {0, + \infty }\right)$ 上单调递减, 由 $f\left( 3\right)  = 0$ ,得 $g\left( 3\right)  = 0$ ,不等式 $\left( {x - 2}\right) f\left( x\right)  < 0 \Longleftrightarrow  \left( {x - 2}\right)  \cdot  \dfrac{f\left( x\right) }{x} < 0 \Longleftrightarrow  \left( {x - 2}\right) g\left( x\right)  < 0$ , 则 $\left\{  \begin{array}{l} x - 2 < 0 \\  g\left( x\right)  > 0 \end{array}\right.$ 或 $\left\{  \begin{array}{l} x - 2 > 0 \\  g\left( x\right)  < 0 \end{array}\right.$ ,即 $\left\{  \begin{matrix} x < 2 \\  0 < x < 3 \end{matrix}\right.$ 或 $\left\{  \begin{array}{l} x > 2 \\  x > 3 \end{array}\right.$ ,解得 $0 < x < 2$ 或 $x > 3$ , 所以不等式 $\left( {x - 2}\right) f\left( x\right)  < 0$ 的解集为 $\left( {0,2}\right)  \cup  \left( {3, + \infty }\right)$ .
        \end{jieda}
        
        \item 已知偶函数 $f\left( x\right)$ 的定义域是 $\left( {-\infty ,0}\right)  \cup  \left( {0, + \infty }\right)$ ,其导函数为 $f\left( x\right)$ ,对定义域内的任意 $x$ ,都有 ${2f}\left( x\right)  + x{f}^{\prime }\left( x\right)  > 2$ 成立,则不等式 ${x}^{2}f\left( x\right)  - {4f}\left( 2\right)  < {x}^{2} - 4$ 的解集为\blkda[4]{$\left( {-2,0}\right)  \cup  \left( {0,2}\right)$}.
        \begin{jieda}
            设$F(x) = x^2f(x)-x^2$，则$F'(x) = 2xf(x)+x^2f'(x)-2x = x[{2f}\left( x\right)  + x{f}^{\prime }\left( x\right)-2]$，故$F(x)$为偶函数，在$(-\infty, 0)$严减，在$(0, +\infty)$严增

            ${x}^{2}f\left( x\right)  - {4f}\left( 2\right)  < {x}^{2} - 4 \Longleftrightarrow F(x) < F(2)$，故解集为$\left( {-2,0}\right)  \cup  \left( {0,2}\right)$
        \end{jieda}
        \item 函数 $f\left( x\right)$ 的导函数 ${f}^{\prime }\left( x\right)$ 满足 ${2f}\left( x\right)  + {f}^{\prime }\left( x\right)  > 2$ ,且 $f\left( 1\right)  = {2025}$ ,则不等式 $f\left( x\right)  > 1 + \dfrac{2024}{{{e}}^{{2x} - 2}}$ 的解集是\blkda{$\left( {1, + \infty }\right)$}
        \begin{jieda}
            设$F(x) = f(x)e^{2x-2}-e^{2x-2}$，则$F'(x) = [f'(x)+2f(x)-2]e^{2x-2}$，故$F'(x) > 0$，$F(1) = 2024$

            $f\left( x\right)  > 1 + \dfrac{2024}{{{e}}^{{2x} - 2}} \Longleftrightarrow F(x) > 2024$，故解集为$\left( {1, + \infty }\right)$

        \end{jieda}
        
        \item 设函数 $f\left( x\right)  = \dfrac{{x}^{2}}{2} - {2a}\ln x + \left( {a - 2}\right) x,a \in  R$ ,若果存在实数 $a$ ,对于不相等正数 ${x}_{1},{x}_{2}$ ,不等式 $\dfrac{f\left( {x}_{1}\right)  - f\left( {x}_{2}\right) }{{x}_{1} - {x}_{2}} > a$ 恒成立,则实数 $a$ 的取值范围为\blkda{$\left(-\infty, -\dfrac{1}{2}\right]$}
        \begin{jieda}
            不妨设$x_1 > x_2$
            
            $\dfrac{f\left( {x}_{1}\right)  - f\left( {x}_{2}\right) }{{x}_{1} - {x}_{2}} > a \Longleftrightarrow f(x_1)-ax_1 > f(x_2)-ax_2$，即$F(x) = f(x)-ax$在$(0, +\infty)$为增函数，故$F'(x) = x-\dfrac{2a}{x}-2 \geq  0$在$x > 0$恒成立，故$x^2-2x \geq 2a$在$x > 0$恒成立，解得$a \in \left(-\infty, -\dfrac{1}{2}\right]$
        \end{jieda}
        
        \item 对于 $x \in  \left( {0, + \infty }\right)$ ,不等式 ${{e}}^{x} - \ln \left( {mx}\right)  + \left( {1 - m}\right) x \geq  0$ 恒成立，则实数 $m$ 的取值范围为\blkda{$0 < m \leq  {e}$}

        \begin{jieda}
            已知 $x \in  \left( {0, + \infty }\right)$ ,由 ${{e}}^{x} - \ln \left( {mx}\right)  + \left( {1 - m}\right) x \geq  0$ 得, ${{e}}^{x} + x \geq  {{e}}^{\ln \left( {mx}\right) } + \ln \left( {mx}\right)$ ,
            
            构造函数 $f\left( x\right)  = {{e}}^{x} + x$ ,则 $f\left( x\right)$ 是 $\mathbf{R}$ 上的增函数,则由 $f\left( x\right)  \geq  f\left( {\ln \left( {mx}\right) }\right)$ 得: $x \geq  \ln \left( {mx}\right)$，即 $m \leq  \dfrac{{{e}}^{x}}{x}$
            
            令 $g\left( x\right)  = \dfrac{{{e}}^{x}}{x},x \in  \left( {0, + \infty }\right) ,{g}^{\prime }\left( x\right)  = \dfrac{\left( {x - 1}\right) {{e}}^{x}}{{x}^{2}}$，当 $x \in  \left( {0,1}\right)$ , ${g}^{\prime }\left( x\right)  < 0$ ,则 $g\left( x\right)$ 单调递减，当 $x \in  \left( {1, + \infty }\right)$ , ${g}^{\prime }\left( x\right)  > 0$ ,则 $g\left( x\right)$ 单调递增
            
            故$g{\left( x\right) }_{\min } = g\left( 1\right)  = {e}$ ,则 $m \leq  {e}$ ,又 $m > 0$ ,则 $0 < m \leq  {e}$ .
        \end{jieda}

        
        \item 已知函数 $f\left( x\right)  = \dfrac{\ln x}{x} - x + 1$ .若对于任意的 $\left( {0, + \infty }\right) ,f\left( x\right)  +x+ \dfrac{1}{x} \leq  a{e}^{x}$ 恒成立,求实数 $a$ 的最小值.
        \begin{jieda}
            $\dfrac{\ln x}{x} - x + x + 1 + \dfrac{1}{x} \leq ae^x$在$\left( {0, + \infty }\right)$ 恒成立

            即$a \geq \dfrac{lnx+x+1}{xe^x}$， 令$t = xe^x$，则$lnt = x + lnx$，故有$a \geq \dfrac{lnt+1}{t}$在$t \in (0, +\infty)$恒成立.

            \ps 若为填空题，此时直曲分离，数形结合可出答案

            设$g(t) = \dfrac{lnt+1}{t}$，$g'(t) = \dfrac{-lnt}{t^2}$，故$t \in (0, 1)$时$g'(t) > 0$，$g(t)$严增，$t \in (1, +\infty)$时$g'(t) < 0$，$g(t)$严减，因此$a \geq g(1) = 1$，所以$a$ 的最小值为$1$
        \end{jieda}
\end{enumerate}
\subsection*{计数原理}
\begin{enumerate}
    \item \begin{enumerate}
        \item 解方程 ${{C}}_{12}^{{x}^{2} - {2x}} = {{C}}_{12}^{{3x} + 6}$
        \item 解不等式 ${C}_{m}^{m - 4} > {C}_{m - 1}^{5} + {C}_{m - 1}^{6}$
    \end{enumerate}
    \begin{jieda}
        \begin{enumerate}
            \item $x = -1, 2$
            \item $m = 7,8,9$
        \end{enumerate}
    \end{jieda}
    
    \item \begin{enumerate}
        \item 计算: ${C}_{3}^{3} + {C}_{4}^{3} + {C}_{5}^{3} + \cdots  + {C}_{98}^{3} + {C}_{99}^{3}$
        \item 证明: ${C}_{m}^{m} + {C}_{m + 1}^{m} + {C}_{m + 2}^{m} + \cdots  + {C}_{m + k - 2}^{m} + {C}_{m + k - 1}^{m} = {C}_{m + k}^{m + 1}$ .
    \end{enumerate}
    \begin{jieda}
        \begin{enumerate}
            \item ${C}_{3}^{3} + {C}_{4}^{3} + {C}_{5}^{3} + \cdots  + {C}_{98}^{3} + {C}_{99}^{3} = {C}_{4}^{4} + {C}_{4}^{3} + {C}_{5}^{3} + \cdots  + {C}_{98}^{3} + {C}_{99}^{3} = {C}_{5}^{4} + {C}_{5}^{3} + \cdots  + {C}_{98}^{3} + {C}_{99}^{3} = \cdots = {C}_{100}^{4} = 3921225$
            \item ${C}_{m}^{m} + {C}_{m + 1}^{m} + {C}_{m + 2}^{m} + \cdots  + {C}_{m + k - 2}^{m} + {C}_{m + k - 1}^{m} = {C}_{m+1}^{m+1} + {C}_{m + 1}^{m} + {C}_{m + 2}^{m} + \cdots  + {C}_{m + k - 2}^{m} + {C}_{m + k - 1}^{m} = {C}_{m + 2}^{m+1} + {C}_{m + 2}^{m} + \cdots  + {C}_{m + k - 2}^{m} + {C}_{m + k - 1}^{m} = \cdots = {C}_{m + k}^{m + 1}$
        \end{enumerate}
    \end{jieda}
    \item 书架上原来摆放着 $6$ 本书，现在要插入 $3$ 本不同的书，则不同的插法为\dotfill{(\kh{C})}
\xx{$P_{7}^{3}$}{$P_{8}^{3}$}{$P_{9}^{3}$}{$2P_{3}^{3}$}


\item 有 $4$ 名男生、$5$ 名女生，全体排成一行，问下列情形中各有多少种不同的排法？
\begin{enumerate*}
\item 甲不在中间也不在两端;
\item 甲、乙两人必须排在两端;
\item 男、女生分别排在一起;
\item 男女相间隔;
\item 甲在乙的右边，甲、乙可以不相邻;
\item 甲、乙、丙三人从左到右顺序保持一定顺序.
\end{enumerate*}

\begin{jieda}
    \begin{enumerate*}
        \item 241920
        \item 10080
        \item 5760
        \item 2880
        \item 181440
        \item 60480
    \end{enumerate*}
\end{jieda}

\item 由 $1$ 到 $9$ 这九个数字中每次选出 $5$ 个组成无重复数字的 $5$ 位数.
\begin{enumerate}
\item 其中奇数位置上只能是奇数，问有多少个这样的 $5$ 位数？
\item 其中奇数只能在奇数位置上，又有多少个这样的 $5$ 位数？
\end{enumerate}
\begin{jieda}
    \begin{enumerate*}
        \item 1800
        \item 2520
    \end{enumerate*}
\end{jieda}

\item 某次文艺晚会上计划演出 $7$ 个节目，其中 $2$ 个歌曲节目，$3$ 个舞蹈节目，$2$ 个小品节目，需要制作节目单：
\begin{enumerate}
\item 唱歌节目排在两头，有多少种排法？
\item 唱歌节目相邻，舞蹈节目相邻，两个小品节目不相邻，有多少种排法？
\item 三个舞蹈节目出场顺序固定，有多少种排法？
\end{enumerate}

\begin{jieda}
    \begin{enumerate*}
        \item 240
        \item 144
        \item 840
    \end{enumerate*}
\end{jieda}

\item 5名乒乓球队员中,有2名老队员和3名新队员,现从中选出3名队员排成1、2、3号参加团体比赛,则入选的3名队员中至少有一名老队员,且1、2号中至少有1名新队员的排法有\blkda{48}种.
\begin{jieda}
    \begin{dingautolist}{172}
        \item 两新一老：$[C_3^2C_2^1] \times [P_3^3] = 36$
        \item 两老一新：$[C_3^1C_2^2] \times [C_2^1P_2^2] = 12$
    \end{dingautolist}
\end{jieda}

\item 2位男生和3位女生共5位同学站成一排,若男生甲不站两端,3位女生中有且只有2位女生相邻,则不同排法的种数是\blkda{48}种.
\begin{jieda}
    \begin{dingautolist}{172}
        \item 将女生分成$2+1$两组：$C_3^1$种方案
        \item 2个女生内部顺序：$P_2^2$
        \item 两个男生排好：$P_2^2$
        \item 两组女生插空
        
        要满足甲不在两端，甲旁边端侧的空必须占位，所以考虑该位置从两组女生中选一组，$C_2^1$，剩下一组女生从余下两个空位中选择一个$C_2^1$
    \end{dingautolist}
    $C_3^1 \times P_2^2 \times P_2^2 \times C_2^1 \times C_2^1 = 48$
\end{jieda}


\item 有4名男生,5名女生.
\begin{enumerate}
    \item 从中选出5名代表,有多少种选法?
    \item 从中选出5名代表,男生2名,女生3名且某女生必须在内有多少种选法?
    \item 从中选出5名代表,男生不少于2名,有多少种选法?
    \item 分成三个小组,每组依次有4、3、2人有多少种分组方法?
\end{enumerate}
\begin{jieda}
    \begin{enumerate}
        \item ${C}_{9}^{5} = {126}$
        \item $ {C}_{4}^{2} \times  1 \times  {C}_{4}^{2} = {36}$
        \item $\left( {{C}_{4}^{2} \times  {C}_{5}^{3}}\right)  + \left( {{C}_{4}^{3} \times  {C}_{5}^{2}}\right)  + \left( {{C}_{4}^{4} \times  {C}_{5}^{1}}\right)  = {105}$
        \item ${C}_{9}^{4} \times  {C}_{5}^{3} \times  {C}_{2}^{2} = {1260}$
    \end{enumerate}
\end{jieda}

\item 四个小球放入编号为1、2、3、4四个盒子中,依下列条件各有多少种放法.
\begin{enumerate*}
    \item 四个小球不同,每个盒子各放一个;
    \item 四个小球相同,每个盒子各放一个;
    \item 四个小球不同,四个盒子恰有一个空着;
    \item 四个小球相同,恰有一个盒子空着.
\end{enumerate*}
\begin{jieda}
    \begin{enumerate}
        \item $P_4^4 = 24$
        \item $C_4^4 = 1$
        \item 先选出一个盒子空着$C_4^1 = 4$，再对小球分组$\dfrac{C_4^2C_2^2C_1^1}{P_2^2} = 6$，再进行分配$P_3^3 = 4$，故$4 \times 6 \times 4 = 144$种
        \item 先选出一个盒子空着$C_4^1 = 4$，再对小球分组，只有$(2,1,1)$一种，再进行分配，相当于从3个盒子中选择1个放2个球，故$C _3^1 = 3$，故$4 \times 1 \times 3 = 12$种
    \end{enumerate}
\end{jieda}

\item 
\begin{enumerate}
    \item 一条长椅上有7个座位,4个人坐,要求3个空位中恰有2个空位相邻,共有多少种不同坐法?
    \item 一条长椅上有9个座位,3个人坐,若相邻2人之间至少有2个空椅子,共有几种不同的坐法?
\end{enumerate}
\begin{jieda}
    \begin{enumerate}
        \item 将两个相邻空位捆绑，无顺序之别，只有1种，问题转化为6个座位，4个人坐，2个有区别的空位不相邻，采用插空法，$P_4^4 \times P_5^2 = 480$
        \item 问题转化为3个人坐5个座位，$P_5^3 = 60$
    \end{enumerate}
\end{jieda}

\item 有15本不同的书,其中6本是数学书,问以下方法数:
\begin{enumerate*}
    \item 分给甲4本,且都不是数学书;
    \item 平均分为3份;
    \item 平均分给3个不同的人;
    \item 甲分2本,乙分7本,丙分6本;
    \item 1人2本,1人7本,1人6本.
\end{enumerate*}
\begin{jieda}
    \begin{enumerate*}
        \item $C_9^4 = 126$
        \item $\dfrac{C_{15}^5C_{10}^5C_5^5}{P_3^3} = 126126$
        \item $\dfrac{C_{15}^5C_{10}^5C_5^5}{P_3^3} \times P_3^3 = 756756$
        \item $C_{15}^2C_{13}^7C_6^6 = 180180$
        \item $C_{15}^2C_{13}^7C_6^6P_3^3 = 1081080$
    \end{enumerate*}
\end{jieda}

\item 
\begin{enumerate}
    \item 求$x_{1}+x_{2}+x_{3}+x_{4}=13$的正整数解的组数;
    \item 求$x_{1}+x_{2}+x_{3}+x_{4}=13$的自然数解的组数;
    \item 求$x_{1}+x_{2}+x_{3}+x_{4}\le13$的正整数解的组数.
\end{enumerate}
\begin{jieda}
    \begin{enumerate}
        \item 看做$13$个相同的小球排成1排，在中间12个空中选择$3$个插入隔板，$C_{12}^3 = 220$
        \item 转化为$x_{1}+x_{2}+x_{3}+x_{4}=17$的正整数解的组数，$C_{16}^3 = 560$
        \item $C_3^3 + C_4^3 + \cdots + C_{12}^3 = C_{13}^4 = 715$
    \end{enumerate}
\end{jieda}

\item 将甲、乙、丙、丁四名学生分到三个不同的班,每个班至少分到一名学生,且甲、乙两名学生不能分到同一个班,则不同分法的种数为\blkda{30}种.
\begin{jieda}
    先分组，共$\dfrac{C_4^2C_2^1C_1^1}{P_2^2} - 1 = 5$种分组方案

    再分配，共$P_3^3 = 6$种分配方案

    $5 \times 6 = 30$.
\end{jieda}

\item 一所医院派出5名医生到3所乡镇卫生院指导,要求每所卫生院院至少派遣一名医生，每个医生只去一个卫生院，则不同的派出方法有\blkda{150}种.
\begin{jieda}
    先分组，共$\dfrac{C_5^3C_2^1C_1^1}{P_2^2} + \dfrac{C_5^2C_3^2C_1^1}{P_2^2} = 25$种分组方案

    再分配，共$P_3^3 = 6$种分配方案

    $25 \times 6 = 150$.
\end{jieda}

\item 四面体的一个顶点为$A$,从其他顶点与棱的中点中取3个点,使它们和点$A$在同一平面上,不同的取法有\blkda{33}种.
\begin{jieda}
    \begin{dingautolist}{172}
        \item 当所取的3点在同一个侧面上时，每个侧面上有 ${{C}}_{5}^{3}$ 种取法，所以共有 $3{{C}}_{5}^{3} = {30}$ 种不同的取法
        \item 当所取的3点不在同一个侧面上时，含顶点 $A$ 的三条棱中,每条棱上有三个点,它们分别与所对的棱的中点共面，共有3种取法
    \end{dingautolist}
    综上,共有 ${30} + 3 = {33}$ 种不同的取法
    \begin{center}
        \includegraphics[width=0.2\linewidth]{bodyxb//src/pailiezuhe2_11.png}
    \end{center}
\end{jieda}

\item 某学校图书室内,有10位同学围着一张圆桌坐成一圈,共有多少种不同的坐法\dotfill{(\kh{B})}
\xx{$P_{10}^{10}$}{$P_{9}^{9}$}{$\dfrac{1}{2}P_{10}^{10}$}{$\dfrac{1}{2}P_{9}^{9}$}


\item 化简: $1-3C_{n}^{1}+3^{2}\cdot C_{n}^{2}-3^{3}\cdot C_{n}^{3}+\cdots+(-1)^{n}3^{n}=$\blkda{$(-2)^n$}

\item 若 $2000<C_{n}^{1}+C_{n}^{2}+\cdots+C_{n}^{n}<3000,$ 则正整数 $n=$\blkda{11}


\item $(1+2x)^{3}(1-x)^{4}$ 展开式中 $x^{2}$ 的系数为\blkda{-6}

\item $\left(\dfrac{\sqrt{x}}{b}+\dfrac{b}{\sqrt[3]{x}}\right)^{18}$ 的展开式中,第\blkda{10}项是中间项,中间项是\blkda{$48620x^{\frac{3}{2}}$}

\item 已知 $(2-\sqrt{3}x)^{100}=a_{0}+a_{1}x+a_{2}x^{2}+\cdots+a_{100}x^{100}$,求下列各式的值: 
\begin{enumerate*}
    \item $a_{0}$ 
    \item $a_{1}+a_{2}+\cdots+a_{100}$ 
    \item $a_{1}+a_{3}+a_{5}+\cdots+a_{99}$ 
    \item $(a_{0}+a_{2}+\cdots+a_{100})^{2}-(a_{1}+a_{3}+\cdots+a_{99})^{2}$ 
    \item $|a_{0}|+|a_{1}|+|a_{2}|+\cdots+|a_{100}|$ 
\end{enumerate*}
\begin{jieda}
    \begin{enumerate}
        \item 令$x = 0$，$a = 2^{100}$
        \item 令$x = 1$，$a_0+a_{1}+a_{2}+\cdots+a_{100} = (2-\sqrt{3})^{100}$，故$a_{1}+a_{2}+\cdots+a_{100} = (2-\sqrt{3})^{100} - 2^{100}$
        \item 令$x = -1$，$a_0-a_{1}+a_{2}+\cdots+(-1)^{99}a_{99}+a_{100} = (2+\sqrt{3})^{100}$
        与$a_0+a_{1}+a_{2}+\cdots+a_{100} = (2-\sqrt{3})^{100}$两式相减除以2可得$a_{1}+a_{3}+a_{5}+\cdots+a_{99} = \dfrac{(2-\sqrt{3})^{100}-(2+\sqrt{3})^{100}}{2}$
        \item $(a_{0}+a_{2}+\cdots+a_{100})^{2}-(a_{1}+a_{3}+\cdots+a_{99})^{2} = (a_0+a_{1}+a_{2}+\cdots+a_{100})(a_0-a_{1}+a_{2}+\cdots+(-1)^{99}a_{99}+a_{100}) = (2-\sqrt{3})^{100}(2+\sqrt{3})^{100} = 1$
        \item $|a_{0}|+|a_{1}|+|a_{2}|+\cdots+|a_{100}|$  = $a_0-a_{1}+a_{2}+\cdots+(-1)^{99}a_{99}+a_{100} = (2+\sqrt{3})^{100}$
    \end{enumerate}
\end{jieda}

\item 已知 $(x^{\frac{2}{3}}+3x^{2})^{n}$ 的展开式各项系数和比它的二项式系数和大992. 
\begin{enumerate}
    \item 求展开式中二项式系数最大的项; 
    \item 求展开式中系数最大的项. 
\end{enumerate}
\begin{jieda}
    令 $x = 1$ 得展开式各项系数和为 ${4}^{n}$ ,二项式系数为 ${C}_{n}^{0} + {C}_{n}^{1} + \ldots \ldots  + {C}_{n}^{n} = {2}^{n}$，由题意得： ${4}^{n} - {2}^{n} = {992}$ ,解得 $n = 5$
    \begin{enumerate}
        \item 展开式共6项,二项式系数最大项为第三、四项
        ${T}_{3} = C_5^2x^29x^4 = 90x^6$，$T_4 = C_5^3x^\frac{4}{3}27x^6 = 270x^\frac{22}{3}$
        \item 设第 $r + 1$ 项的系数为 ${t}_{r + 1}$ ,则 ${t}_{r + 1} = {C}_{5}^{r}{3}^{r}$
        
        $\dfrac{{t}_{r + 2}}{{t}_{r + 1}} = \dfrac{C_5^{r+1}3^{r+1}}{{C}_{5}^{r}{3}^{r}} = 3 \times \dfrac{5-r}{r+1}$，令$3 \times \dfrac{5-r}{r+1} = 1$，即$15-3r = r+1$，$r = 3.5$，因此$t_5$最大，故展开式系数最大项为： ${T}_{5} = {405}x^\frac{26}{3}$
    \end{enumerate}
\end{jieda}
\end{enumerate}
\subsection*{概率}
\begin{enumerate}
    \item 下列说法不正确的是\dotfill{(\kh{ABD})}
    \xx{频率就是概率}{任何事件的概率都是在$(0,1)$之间}{概率是客观存在的，与试验次数无关}{概率是随机的，与试验次数有关}
    
    \item 下列说法正确的是\dotfill{(\kh{D})}
    \xx{某人打靶，射击$10$次，中靶$7$次，则此人中靶的概率为$0.7$}{一位同学做抛硬币试验，抛$6$次，一定有$3$次“正面朝上”}{某地发行一种彩票，回报率为$47\%$，若有人花$100$元钱买此种彩票，则一定会有$47$元的回报}{用某种药物对患有胃溃疡的$500$名病人进行治疗，结果有$380$人有明显的疗效，现在胃溃疡病人服用此药，则可估计有明显疗效的概率约为$0.76$}
    
    \item 甲、乙两人玩猜数字游戏,先由甲心中想一个数字,记为$a$,再由乙猜甲刚才所想的数字,把乙猜的数字记为$b$,其中 $a,b\in\{1,2,3,4,5,6\}$,若$|a-b|\le1$,就称“甲、乙心有灵犀”,现任意找两人玩这个游戏,则他们“心有灵犀”的概率为\blkda{$\dfrac{4}{9}$}.
    \item 吸烟有害健康,小明为了帮助爸爸戒烟,在爸爸包里放一个小盒子,里面随机摆放三支香烟和三支跟香烟外形完全一样的“戒烟口香糖”,并且和爸爸约定,每次想吸烟时,从盒子里任取一支,若取到口香糖则吃一支口香糖,不吸烟;若取到香烟,则吸一支烟,不吃口香糖,假设每次香烟和口香糖被取到的可能性相同,则“口香糖吃完时还剩2支香烟”的概率为\blkda{$\dfrac{3}{20}$}.
        
    \item 从4名男生和2名女生中任选3人参加演讲比赛.
    \begin{enumerate}
    \item 求所选3人都是男生的概率;
    \item 求所选3人中恰有1名女生的概率;
    \item 求所选3人中至少有1名女生的概率.
    \end{enumerate}
    \begin{jieda}
        \begin{enumerate*}
            \item $\dfrac{1}{5}$
            \item $\dfrac{3}{5}$
            \item $\dfrac{4}{5}$
        \end{enumerate*}
    \end{jieda}
    
    
    \item 某人有5把钥匙,但忘记了开房门的是哪一把,于是,他逐把不重复地试开,问:
    \begin{enumerate}
    \item 恰好第三次打开房门所的概率是多少?
    \item 三次内打开的概率是多少?
    \item 如果5把内有2把房门钥匙,那么三次内打开的概率是多少?
    \end{enumerate}
    \begin{jieda}
        \begin{enumerate*}
            \item $\dfrac{1}{5}$
            \item $\dfrac{3}{5}$
            \item $\dfrac{9}{10}$
        \end{enumerate*}
    \end{jieda}
        
    \item 用3种不同的颜色给甲、乙两个小球随机涂色,每个小球只涂一种颜色,则两个小球颜色不同的概率为\dotfill{(\kh{C})}
    \xx{$\dfrac{1}{3}$}{$\dfrac{1}{2}$}{$\dfrac{2}{3}$}{$\dfrac{5}{8}$}
    
    \item 袋中有五张卡片,其中红色卡片三张,标号分别为1、2、3;蓝色卡片两张,标号分别为4、5.从这五张卡片中任取两张,如果5张卡片被抽取的机会相等,分别计算下列事件的概率:
    \begin{enumerate}
    \item 事件$A=$“这两张卡片颜色不同”;
    \item 事件$B=$“这两张卡片标号之和小于7”;
    \item 事件$C=$“这两张卡片颜色不同且标号之和小于7”.
    \end{enumerate}
    
    \begin{jieda}
        \begin{enumerate*}
            \item $\dfrac{3}{5}$
            \item $\dfrac{3}{5}$
            \item $\dfrac{3}{10}$
        \end{enumerate*}
    \end{jieda}
    
    \item 现从 A, B, C, D, E五人中任选三人参加一个重要会议,五人被选中的机会相等.求:
    \begin{enumerate*}
    \item A被选中的概率;
    \item A和B同时被选中的概率;
    \item A或B被选中的概率.
    \end{enumerate*}
    
    \begin{jieda}
        \begin{enumerate*}
            \item $\dfrac{3}{5}$
            \item $\dfrac{3}{10}$
            \item $\dfrac{9}{10}$
        \end{enumerate*}
    \end{jieda}
        
    \item 甲、乙、丙三台机床各自独立地加工同一种零件,已知甲机床加工的零件是一等品而乙机床加工的零件不是一等品的概率为 $\dfrac{1}{4}$, 乙机床加工的零件是一等品而丙机床加工的零件不是一等品的概率为 $\dfrac{1}{12}$, 甲、丙两台机床加工的零件都是一等品的概率为 $\dfrac{2}{9}$.
    \begin{enumerate}
    \item 分别求甲、乙、丙三台机床各自加工零件是一等品的概率;
    \item 从甲、乙、丙加工的零件中各取一个检验,求至少有一个一等品的概率.
    \end{enumerate}
    \begin{jieda}
        \begin{enumerate}
            \item 由事件的独立性，$P(A)=\dfrac{1}{3}$,$P(B)=\dfrac{1}{4}$,$P(C)=\dfrac{2}{3}$
            \item 由对立事件和独立性，$\dfrac{5}{6}$
        \end{enumerate}
    \end{jieda}

    
    \item 棱长为1的正四面体A-BCD,有一小虫从顶点A处开始按以下规则爬行: 在每一顶点处以同样的概率选择通过这个顶点的3条棱之一,并一直爬到这条棱的尽头,记小虫爬了$n$米后重新回到点A的概率为$P_n$.
    \begin{enumerate*}
    \item 求$P_1$和$P_2$的值;
    \item 探寻$P_n$与$P_{n-1}$的关系;
    \item 求$P_n$的表达式.
    \end{enumerate*}
    
    \begin{jieda}
        \begin{enumerate*}
            \item $P_1=0$,$P_2=\dfrac{1}{3}$
            \item $P_n=\dfrac{1}{3}(1-P_{n-1})$
            \item $P_{n}=\dfrac{1}{4}+\dfrac{3}{4}(-\dfrac{1}{3})^n$
        \end{enumerate*}
    \end{jieda}

        
    \item 现有$6$个分别标有数字$1,2,3,4,5,6$的相同小球，从中有放回地随机抽取两次，每次取$1$个球，记事件甲：第一次取出的球的数字是$3$，事件乙：第二次取出的球的数字是$6$，事件丙：两次取出的球的数字之和是$8$，事件丁：两次取出的球的数字之差的绝对值是$3$，则正确的有\blkda{(2)(4)}
    
    \begin{enumerate*}
        \item 甲与丙相互独立
        \item 甲与丁相互独立
        \item 乙与丙相互对立
        \item 丙与丁互斥
    \end{enumerate*}
    
        
    \item 甲、乙二人进行一场游戏比赛，且比赛中不存在平局，先赢三局者获胜，并可以获得$800$元奖金，已知甲、乙二人在每局比赛中获胜的可能性均相同，已知当甲连赢两局，乙一局未赢时，因某种特殊情况需要终止比赛，现将$800$元奖金按两人各自最终获胜的可能性的比例进行分配，则甲应该分得\blkda{700}元

    \item 假设每一架飞机的每一个引擎在飞行中出现故障概率均为$1-p$,且各引擎是否有故障是独立的,已知4引擎飞机中至少有3个引擎飞机正常运行,飞机就可成功飞行;2引擎飞机要2个引擎全部正常运行,飞机才可成功飞行,要使4引擎飞机比2引擎飞机更安全,则 $p$ 的取值范围是\blkda{$(\dfrac{1}{3},1)$}.
    
    \item $4$粒种子种在甲坑内，每粒种子发芽的概率为$\dfrac{1}{2}$，若坑内至少有$2$粒种子发芽，则不需要补种；否则需要补种，则甲坑不需要补种的概率为\blkda{$\dfrac{11}{16}$}
    
    \item 某自助银行有A，B，C，D四台ATM，在某一时刻这四台ATM被占用的概率分别为$\dfrac{1}{3},\dfrac{1}{2},\dfrac{1}{2},\dfrac{2}{5}$
    \begin{enumerate}
    \item 若某客户只能使用四台ATM中的A或B，则该客户需要等待的概率为\blkda{$\dfrac{1}{6}$}
    \item 某客户使用ATM取款时，恰好有两台ATM被占用的概率为\blkda{$\dfrac{11}{30}$}
    \end{enumerate}

    \item 甲乙二人争夺一场围棋比赛的冠军，若比赛为“三局两胜”制（无平局），甲在每局比赛中获胜的概率均为 $\dfrac{2}{3}$ 且各局比赛结果相互独立，则在甲获得冠军的条件下，比赛进行了三局的概率为\blkda{$\dfrac{2}{5}$}

    \item 已知随机事件 $A$，$B$，若 $P(B|A)=\dfrac{3}{5},P(\overline{A}|B)=\dfrac{4}{7},P(B)=\dfrac{7}{15}$，则$P(A)=$\blkda{$\dfrac{1}{3}$}
    
    \item 设$A$，$B$是两个随机事件，已知 $P(A)=\dfrac{3}{5}$， $P(B)=\dfrac{2}{3}$， $P(\overline{A}|B)=\dfrac{2}{5}$ 记 $C=A\cup B$，则$P(A|C)=$\blkda{$\dfrac{9}{13}$}
    
    \item 一袋中有$6$个黑球，$4$个白球。
    \begin{enumerate}
    \item 依次取出$3$个球，不放回，已知第一次取出的是白球，求第三次取到黑球的概率；
    \item 有放回地依次取出$3$球，已知第一次取的是白球，求第三次取到黑球的概率。
    \end{enumerate}
    \begin{jieda}
        \begin{enumerate}
            \item 设事件 $A$ 为 “第一次取到白球”,事件 $B$ 为 “第二次取
    到白球”,事件 $C$ 为 “第三次取到白球”.
    则在事件 $A$ 发生的条件下,袋中只剩 6 个黑球和 3 个白球,则 $P\left( {\bar{C} \mid  A}\right)  = \dfrac{n\left( {A\bar{C}}\right) }{n\left( A\right) } = \dfrac{{{C}}_{4}^{1}\left( {{{C}}_{3}^{1}{{C}}_{6}^{1} + {{A}}_{6}^{2}}\right) }{{{C}}_{4}^{1} \cdot  {{A}}_{9}^{2}} = \dfrac{2}{3}$ .
            \item 因为每次取之前袋中球的情况不变,
    所以 $n$ 次取球的结果互不影响,所以 $P\left( \bar{C}\right)  = \dfrac{6}{10} = \dfrac{3}{5}$ .
        \end{enumerate}
    \end{jieda}

    
    \item 一台机床有 $\dfrac{1}{3}$ 的时间加工零件A,其余时间加工零件B,加工A时,停机的概率是 $\dfrac{3}{10}$, 加工B时,停机的概率是 $\dfrac{2}{5}$, 则这台机床停机的概率为\dotfill{(\kh{A})}
    \xx{$\dfrac{11}{30}$}{$\dfrac{7}{30}$}{$\dfrac{7}{10}$}{$\dfrac{1}{10}$}
    \begin{jieda}
        全概率公式
    \end{jieda}
    
    \item 有甲、乙两个袋子，甲袋中有$2$个白球，$1$个红球，乙袋中有$2$个红球，一个白球，这$6$个球手感上不可区别，今从甲袋中任取一球放入乙袋，搅匀后再从乙袋中任取一球，此球是红球的概率为\blkda{$\dfrac{7}{12}$}，若已知取到一个红球，则从甲袋放入乙袋的是白球的概率为\blkda{$\dfrac{4}{7}$}

    \item 已知甲盒中有2个红球,1个蓝球,乙盒中有1个红球,2个蓝球,从甲乙两个盒中各取1球放入原来为空的丙盒中,现从甲盒中取1个球,记红球的个数为 $\xi_{1}$ ,从乙盒中取1个球,记红球的个数为 $\xi_{2}$ ,从丙盒中取1个球,记红球的个数为 $\xi_{3}$,则下列说法正确的是\dotfill{(\kh{C})}
    \xx{$E[\xi_{1}]>E[\xi_{3}]>E[\xi_{2}],D[\xi_{1}]=D[\xi_{2}]>D[\xi_{3}]$}{$E[\xi_{1}]<E[\xi_{3}]<E[\xi_{2}],D[\xi_{1}]=D[\xi_{2}]>D[\xi_{3}]$}{$E[\xi_{1}]>E[\xi_{3}]>E[\xi_{2}],D[\xi_{1}]=D[\xi_{2}]<D[\xi_{3}]$}{$E[\xi_{1}]<E[\xi_{3}]<E[\xi_{2}],D[\xi_{1}]=D[\xi_{2}]<D[\xi_{3}]$}
    \begin{jieda}
        随机变量 ${\xi }_{1}$ 可取值0,1,其中 $P\left( {{\xi }_{1} = 0}\right)  = \dfrac{2}{3} \times  \dfrac{1}{2} = \dfrac{1}{3},P\left( {{\xi }_{1} = 1}\right)  = \dfrac{1}{3} \times  1 + \dfrac{2}{3} \times  \dfrac{1}{2} = \dfrac{2}{3}$ ,
    故 $E\left( {\xi }_{1}\right)  = \dfrac{2}{3},D\left( {\xi }_{1}\right)  = \dfrac{2}{3} - \dfrac{4}{9} = \dfrac{2}{9}$ .
    随机变量 ${\xi }_{2}$ 可取值0,1,其中 $P\left( {{\xi }_{2} = 0}\right)  = \dfrac{2}{3} \times  \dfrac{1}{2} + \dfrac{1}{3} \times  1 = \dfrac{2}{3}$ ,
    $P\left( {{\xi }_{2} = 1}\right)  = \dfrac{2}{3} \times  \dfrac{1}{2} = \dfrac{1}{3},$
    故 $E\left( {\xi }_{2}\right)  = \dfrac{1}{3},D\left( {\xi }_{2}\right)  = \dfrac{1}{3} - \dfrac{1}{9} = \dfrac{2}{9}$ .
    随机变量 ${\xi }_{3}$ 可取值 0,1,当 ${\xi }_{3} = 0$ 时,丙盒中无红球或有一个红球,
    无红球的概率为 $\dfrac{1}{3} \times  \dfrac{2}{3}$ ,有一个红球的概率为 $\dfrac{2 \times  2}{9} + \dfrac{1}{9}$ ,
    故 $P\left( {{\xi }_{3} = 0}\right)  = \dfrac{1}{3} \times  \dfrac{2}{3} \times  1 + \left( {\dfrac{2 \times  2}{9} + \dfrac{1}{9}}\right)  \times  \dfrac{1}{2} = \dfrac{1}{2},P\left( {{\xi }_{3} = 1}\right)  = 1$
    $- \dfrac{1}{2} = \dfrac{1}{2},$
    故 $E\left( {\xi }_{3}\right)  = \dfrac{1}{2},D\left( {\xi }_{3}\right)  = \dfrac{1}{2} - \dfrac{1}{4} = \dfrac{1}{4}$ .
    综上所述, $E\left( {\xi }_{1}\right)  > E\left( {\xi }_{3}\right)  > E\left( {\xi }_{2}\right) ,D\left( {\xi }_{1}\right)  = D\left( {\xi }_{2}\right)  < D\left( {\xi }_{3}\right)$ ,故选 C.
    \end{jieda}

    \item 已知离散型随机变量$X$的分布列为: $\begin{pmatrix}X&0&1&2\\ P&a&b&c\end{pmatrix}$ 且$E[X]=\dfrac{3}{2}$ 则当$a^{2}+b^{2}+c^{2}$ 取最小值时,方差 $D[X]=$\blkda{$\dfrac{5}{12}$}

    \item 某超市计划按月订购一种酸奶,每天进货量相同,进货成本每瓶4元,售价每瓶6元,未售出的酸奶降价处理,以每瓶2元的价格当天全部处理完,根据往年销售经验,每天需求量与当天最高气温(单位:C)有关,如果最高气温不低于25,需求量为500瓶;如果最高气温位于区间$[20,25)$,需求量为300瓶;如果最高气温低于20,需求量为200瓶,为了确定六月份的订购计划,统计了前三年六月份各天的最高气温数据,得下面的频数分布表:
    \begin{center}
    \begin{tabular}{|c|c|c|c|c|c|c|}
    \hline
    最高气温 & $[10, 15)$ & $[15, 20)$ & $[20, 25)$ & $[25, 30)$ & $[30, 35)$ & $[35, 40)$ \\
    \hline
    天数 & $25$ & $16$ & $2$ & $36$ & $7$ & $4$ \\
    \hline
    \end{tabular}
    \end{center}
    以最高气温位于各区间的频率代替最高气温位于该区间的概率.
    \begin{enumerate}
    \item 求六月份这种酸奶一天的需求量$X$(单位:瓶)的分布列;
    \item 设六月份一天销售这种酸奶的利润为$Y$(单位:元),当六月份这种酸奶一天的进货量(单位:瓶)为多少时,$Y$的数学期望达到最大值?
    \end{enumerate}
    
    \begin{jieda}
        \begin{enumerate}
            \item 由题意, $X$ 可能的取值为 200,300,500,由表格数据知$P\left( {X = {200}}\right)  = \dfrac{2 + {16}}{90} = \dfrac{1}{5},P\left( {X = {300}}\right)  = \dfrac{36}{90} = \dfrac{2}{5},$$P\left( {X = {500}}\right)  = \dfrac{{25} + 7 + 4}{90} = \dfrac{2}{5},$所以 $X$ 的分布列为
            $\begin{pmatrix}
             200 & 300& 500\\
             \dfrac{1}{5} & \dfrac{2}{5} & \dfrac{2}{5}
            \end{pmatrix}$
    
            \item 由题意,需求量的最小值为 200,最大值为 500,所以只需考虑 ${200} \leq  n \leq  {500}$ 的情形,当 ${300} \leq  n \leq  {500}$ 时,若最高气温不低于 25,则 $Y =$${6n} - {4n} = {2n}$ ;若最高气温位于 $\lbrack {20},{25})$ ,则 $Y = 2 \times  {300} - 2 \times  (n -$${300}) = {1200} - {2n}$ ;若最高气温低于 20,则 $Y = 2 \times  {200} - 2 \times  \left( {n - {200}}\right)  =$${800} - {2n}$ .故 $E\left( Y\right)  = \dfrac{2}{5} \times  {2n} + \dfrac{2}{5}\left( {{1200} - {2n}}\right)  + \dfrac{1}{5}\left( {{800} - {2n}}\right)$$= {640} - \dfrac{2}{5}n$ ,所以当 $n = {300}$ 时, $E\left( Y\right)$ 最大,为 520.当 ${200} \leq  n < {300}$ 时,若最高气温不低于 20,则 $Y =$${6n} - {4n} = {2n}$ ;若最高气温低于 20,则 $Y = 2 \times  {200} - 2\left( {n - {200}}\right)  =$${800} - {2n}$ .故 $E\left( Y\right)  = \dfrac{4}{5} \times  {2n} + \dfrac{1}{5}\left( {{800} - {2n}}\right)  = {160} + \dfrac{6}{5}n <$${160} + \dfrac{6}{5} \times  {300} = {520}.$综上所述,当 $n = {300}$ 时, $Y$ 的数学期望达到最大值 520.
        \end{enumerate}
    \end{jieda}
        
    \item “石头、剪刀、布”又称“猜丁壳”,是一种流行多年的猜拳游戏,其游戏规则:出拳之前双方齐喊口令,然后在语音刚落时同时出拳,握紧的拳头代表“石头”,食指和中指伸出代表“剪刀”,五指伸开代表“布”.“石头”胜“剪刀”,“剪刀”胜“布”,而“布”又胜过“石头”,若所出的拳相同,则为和局.小军和大明两位同学进行五局三胜制的“石头、剪刀、布”游戏比赛,则小军和大明比赛至第四局小军胜出的概率是\blkda{$\dfrac{2}{27}$}.

    \item 甲、乙两队进行篮球决赛,采取七场四胜制(当一队赢得四场胜利时,该队获胜,决赛结束).根据前期比赛成绩,甲队的主客场安排依次为“主主客客主客主”.设甲队主场取胜的概率为0.6,客场取胜的概率为0.5,且各场比赛结果相互独立,则甲队以 $4:1$ 获胜的概率是\blkda{0.18}.
    
    \item 已知随机变量 $X\sim B(6,0.8)$,若 $P(X=k)$ 最大,则 $D[kX+1]=$\blkda{24}.
    
    \item 某人在$n$次射击中击中目标的次数为$X$,且 $X\sim B(n,0.8)$,记 $P_{k}=P(X=k)$, $k=0,1,2,\dots,n$,若是唯一的最大值,则 $D(X)$ 的值为\blkda{1.28}.

    \item 高尔顿板是英国生物统计学家高尔顿设计用来研究随机现象的模型,在一块木板上钉着若干排相互平行但相互错开的圆柱形小木块,小木块之间留有适当的空隙作为通道,前面挡着一块玻璃,让一个小球从高尔顿板上方的通道口落下,最后掉入高尔顿板下方的某一球槽内,如图所示的高尔顿板有5层小木块,将小球从顶端放入,小球下落的过程中,每次碰到小木块后都等可能地向左或向右落下,最后掉入编号为1、2、$\dots$、5的球槽内.
    
    \begin{center}
    \includegraphics[width=0.1\linewidth]{bodyxb/src/faljfliajlghtg.jpeg}
    \end{center}
    
    \begin{enumerate}
    \item 如图,让一个小球从高尔顿板上方的通道口落下,设小球落入球槽的号码为$X$,求$X$的分布列与数学期望.
    \item 现小禹同学对高尔顿板进行改进,小球在下落的过程中与小木块碰撞时,有 $\dfrac{1}{3}$ 的概率向左, $\dfrac{2}{3}$ 的概率向右滚下,小球共经过4次碰撞后,最后掉入编号为1、2、$\dots$、5的球槽内.将80个小球依次从高尔顿板上方的通道口落下,试问2号球槽中落入多少个小球的概率最大?
    \end{enumerate}
    
    \begin{jieda}
        \begin{enumerate}
            \item 若小球落入球槽的号码为 $X = k,k = 1,2,3,4,5$ ,
            则小球共经过 4 次碰撞,向右 $k - 1$ 次,可得 $P\left( {X = k}\right)  = \dfrac{{{C}}_{4}^{k - 1}}{{2}^{4}} = \dfrac{{{C}}_{4}^{k - 1}}{16}$ ,
            则 $P\left( {X = 1}\right)  = \dfrac{{{C}}_{4}^{0}}{16} = \dfrac{1}{16};P\left( {X = 2}\right)  = \dfrac{{{C}}_{4}^{1}}{16} = \dfrac{1}{4};P\left( {X = 3}\right)  = \dfrac{{{C}}_{4}^{2}}{16} = \dfrac{3}{8};P\left( {X = 4}\right)  = \dfrac{{{C}}_{4}^{3}}{16} = \dfrac{1}{4},P\left( {X = 5}\right)  = \dfrac{{{C}}_{4}^{4}}{16} = \dfrac{1}{16}$ ,
            所以 $X$ 的分布列为
            \begin{pmatrix}
                1& 2& 3& 4& 5\\
                \dfrac{1}{16}&\dfrac{1}{4}&\dfrac{3}{8}&\dfrac{1}{4} & \dfrac{1}{16}
            \end{pmatrix}
            
            $X$ 的期望为 $E\left( X\right)  = 1 \times  \dfrac{1}{16} + 2 \times  \dfrac{1}{4} + 3 \times  \dfrac{3}{8} + 4 \times  \dfrac{1}{4} + 5 \times  \dfrac{1}{16} = 3$ .
            
            \item 若小球落入 2 号球槽,则小球共经过 4 次碰撞,向右 1 次,且每次向右的概率均为 $\dfrac{2}{3}$ ,
            则小球落入 2 号球槽的概率为 $P = {{C}}_{4}^{1} \times  \dfrac{2}{3} \times  {\left( \dfrac{1}{3}\right) }^{3} = \dfrac{8}{81}$ ,
            设 80 个小球落入 2 号球槽的个数为 $Y$ ,则 $Y \sim  B\left( {{80},\dfrac{8}{81}}\right)$
            
            $\dfrac{P(Y=k+1)}{P(Y=k)}=\dfrac{8(80-k)}{73(k+1)}>1$,解得$k<7$;\\
            $\dfrac{P(Y=k+1)}{P(Y=k)}=\dfrac{8(80-k)}{73(k+1)}=1$,解得$k=7$;\\
            $\dfrac{P(Y=k+1)}{P(Y=k)}=\dfrac{8(80-k)}{73(k+1)}<1$,解得$k>7$,根据单调性可得最大项为$k=7\text{or}8$。
        \end{enumerate}
    \end{jieda}

        \item 下列随机事件中的随机变量 $X$ 服从超几何分布的是\dotfill{(\kh{B})}

        \xx{将一枚硬币连抛 3 次，记正面向上的次数为 $X$}
        {从 7 男 3 女共 10 名学生干部中随机选出 5 名学生干部，记选出女生的人数为 $X$}
        { 某射手的射击命中率为 0.8 ，现对目标射击 1 次，记命中的次数为 $X$}
        {盒中有 4 个白球和 3 个黑球，每次从中摸出 1 个球且不放回，记第一次摸出黑球时摸取的次数为 $X$}
        
        
        \item 在 15 个村庄中，有 7 个村庄交通不方便，若用随机变量 $X$ 表示任选 10 个村庄中交通不方便的村庄的个数，则 $X$ 服从参数为 $N$ ， $M$ ， $n$ 的超几何分布，则 $P\left( {X = k}\right)  = \dfrac{{C}_{M}^{k}{C}_{N - M}^{n - k}}{{C}_{N}^{n}}$ ，其参数为\dotfill{(\kh{A})}
        
        \xx{${N} = {15},{M} = 7,{n} = {10}$}
        {${N} = {15},{M} = {10},{n} = 7$}
        {${N} = {22},{M} = {10},{n} = 7$ }
        { ${N} = {22},{M} = 7,{n} = {10}$}  

        
\item 某大学志愿者协会有 10 名同学，成员构成如下表，表中部分数据不清楚，只知道从这 10 名同学中随机选取 1 名同学，该名同学的专业为数学的概率为 $\frac{2}{5}$ .

\begin{center}
    \includegraphics[width=0.3\textwidth]{bodyxb//src/g'apfur]0q83yhfuhoouha.png}
\end{center}

现从这 10 名同学中随机选取 3 名同学参加社会公益活动 (每名同学被选到的可能性相同).
\begin{enumerate}
    \item 求 $m,n$ 的值；
    \item 求选出的 3 名同学恰为专业互不相同的男生的概率；
    \item 设 $\xi$ 为选出的 3 名同学中是女生或专业为数学的人数，求随机变量 $\xi$ 的分布列.
\end{enumerate}

\begin{jieda}
    \begin{enumerate}
        \item 设事件 $A$ 为 “从 10 名同学中随机选取 1 名
同学,该名同学的专业为数学”.
由题意,可知数学专业的同学共有 $\left( {1 + m}\right)$ 名,
则 $P\left( A\right)  = \dfrac{1 + m}{10} = \dfrac{2}{5}$ ,解得 $m = 3$ .
因为 $m + n + 6 = {10}$ ,所以 $n = 1$ .
        \item 设事件 $B$ 为 “选出的 3 名同学恰为专业互不相
同的男生”,则事件 $B$ 包含两种情况: 一是从数学专
业的 3 人中选 1 人,再从剩下的三个专业中选 2 人,
有 ${{C}}_{3}^{1}{{C}}_{3}^{2}$ 种方法,二是选出除数学外的 3 个专业的 3
人,有 1 种选法,
所以 $P\left( B\right)  = \dfrac{{{C}}_{3}^{1}{{C}}_{3}^{2} + 1}{{{C}}_{10}^{3}} = \dfrac{1}{12}$ .

\item 由题意,可知这 10 名同学中是女生或专业为数
学的人数为 $7,\xi$ 的取值范围为 $\{ 0,1,2,3\}$ .
$P\left( {\xi  = 0}\right)  = \dfrac{{{C}}_{3}^{3}}{{{C}}_{10}^{3}} = \dfrac{1}{120},$
$P\left( {\xi  = 1}\right)  = \dfrac{{{C}}_{7}^{1}{{C}}_{3}^{2}}{{{C}}_{10}^{3}} = \dfrac{21}{120} = \dfrac{7}{40},$
$P\left( {\xi  = 2}\right)  = \dfrac{{{C}}_{7}^{2}{{C}}_{3}^{1}}{{{C}}_{10}^{3}} = \dfrac{63}{120} = \dfrac{21}{40},$
$P\left( {\xi  = 3}\right)  = \dfrac{{{C}}_{7}^{3}}{{{C}}_{10}^{3}} = \dfrac{35}{120} = \dfrac{7}{24},$
所以 $\xi$ 的分布列为
$\begin{pmatrix}
    0 &1& 2& 3 \\
    \dfrac{1}{120} & \dfrac{7}{40} & \dfrac{21}{40} & \dfrac{7}{24}
\end{pmatrix}$

    \end{enumerate}

\end{jieda}
    
        \item 已知甲、乙两种水果的质量分别服从正态分布 $N\left( {{\mu }_{1},{\sigma }_{1}^{2}}\right) ,N\left( {{\mu }_{2},{\sigma }_{2}^{2}}\right)$ ,其正态分布的密度曲线如图所示, 则\dotfill{(\kh{A})}
        
        \begin{center}
        \includegraphics[width=0.2\textwidth]{bodyxb/src/bo_d7enkus91nqc73eqra40_0_1069_791_332_208_0.jpg}
        \end{center}
    
        \xx{${\mu }_{1} < {\mu }_{2},{\sigma }_{1} < {\sigma }_{2}$}{${\mu }_{1} < {\mu }_{2},{\sigma }_{1} > {\sigma }_{2}$}{${\mu }_{1} > {\mu }_{2},{\sigma }_{1} < {\sigma }_{2}$}{${\mu }_{1} > {\mu }_{2},{\sigma }_{1} > {\sigma }_{2}$}
            
        \item 在某市2019年1月份的高三质量检测考试中，理科生的数学成绩服从正态分布 $N(98,100)$.已知参加本次考试的全市理科生约有9450 人，如果某学生在这次考试中的数学成绩是108分，那么他的数学成绩大约排在全市第\dotfill{(\kh{A})}
        \xx{1500名}{1700名}{4500名}{8000名}
        \begin{jieda}
            $108 = 98+10 = \mu + \sigma$，$P(\mu - \sigma < x < \mu + \sigma) = 0.6826$，$P(x < \mu + \sigma) = 1-(1-0.6826)\div 2 = 0.8413$，$(1-0.8413) \times 9450 = 1499.7$
        \end{jieda}
\end{enumerate}